\section{Related Work}
\label{sec:related}

\paragraph*{Classical and production UV unwrapping.}
Classical UV unwrapping and production tools optimize geometric objectives: low area and angle distortion, manageable seams, good packing, and stable chart layouts. Foundational parameterization methods include MIPS \cite{hormann2000_mips_efficient_global_parameterization}, angle-based flattening \cite{sheffer2001_abf_parameterization}, LSCM \cite{levy2002_lscm_texture_atlas}, stretch-driven charting and Iso-charts \cite{sander2001_texture_mapping_progressive_meshes,zhou2004_iso_charts}, mean-value parameterization \cite{floater1997_parametrization_smooth_approximation}, local/global ARAP-style parameterization \cite{liu2008_local_global_parameterization}, scalable locally injective mappings \cite{rabinovich2017_scalable_locally_injective}, joint cut/parameterization optimization in OptCuts \cite{li2018_optcuts}, and Boundary First Flattening \cite{sawhney2017_boundary_first_flattening}. Production tools such as xatlas \cite{young2024_xatlas}, Blender Smart UV \cite{blender2026_smart_uv_project}, Maya UV tools \cite{autodesk2026_maya_uv_editor}, and ZBrush UV Master \cite{maxon2026_zbrush_uv_master} operationalize these goals with robust automatic charting, packing, and editing interfaces. In our experiments, xatlas is the strongest production-style baseline on the three main assets, while Blender Smart UV provides a widely used automatic alternative. These tools are valuable precisely because they produce robust atlases with familiar DCC constraints. Our paper does not argue that these objectives are obsolete. Instead, we show that they can be insufficient for PBR deployment: when a part-aware atlas produces residual-heavy PBR baking artifacts, held-out relighting error can indicate where local repair is warranted.

\paragraph*{Texture atlases, packing, and signal metrics.}
Texture atlas work has long connected parameterization to surface signal sampling, including progressive-mesh texture mapping \cite{sander2001_texture_mapping_progressive_meshes}, signal-specialized parameterization \cite{sander2002_signal_specialized_parameterization}, atlas packing refinement \cite{limper2018_box_cutter_atlas_refinement}, atlas-domain processing \cite{prada2018_gradient_domain_texture_atlas}, and texture-atlas compression \cite{luo2023_texture_atlas_compression}. Our evaluation also uses image-space and perceptual metrics with established limitations: PSNR is standard but content-dependent \cite{huynh2008_psnr_validity}, SSIM measures structural similarity \cite{wang2004_ssim}, LPIPS measures learned perceptual similarity \cite{zhang2018_lpips}, and FLIP \cite{andersson2020_flip} is a graphics-oriented alternative when image differences must be evaluated under alternating views. Tangent-space conventions for normal-map baking follow the MikkTSpace specification \cite{mikkelsen2008_mikktspace}, which is the de facto standard in modern game engines and production DCC tools. We therefore report them as deployment diagnostics rather than treating any one metric as a complete atlas-quality proxy.

\paragraph*{Neural and data-driven UV parameterization.}
Recent neural and data-driven parameterization methods target more difficult surfaces, including generated or noisy meshes. PartUV explicitly addresses AI-generated meshes through part-based charting and recursive geometric heuristics \cite{wang2025_partuv_part_based_uv}. FlexPara learns flexible neural parameterizations and chart assignments \cite{zhao2025_flexpara_flexible_neural_surface}. Flatten Anything and ParaPoint are relevant but were not fully reproduced in our environment; their reproduction status is reported in \cref{app:baseline_failure}. These methods primarily reason about parameterization structure. We keep their outputs, or the available production baselines, as initial atlases and ask a downstream question: whether the baked PBR residual justifies a local, deployment-gated repair.

\paragraph*{Capacity allocation and texture-aware UV ideas.}
OT-UVGS is adjacent because it treats UV mapping for Gaussian splatting as a fixed-budget capacity-allocation problem \cite{kim2026_ot_uvgs_revisiting_uv}. That perspective is close to our C3 texel allocator, but the optimization target differs. We allocate mesh-chart texels using held-out PBR residual, visibility, channel frequency, and mip leakage under a single DCC-ready atlas. OT-UVGS allocates UV capacity for Gaussian primitives, whereas our repair is chart-local and validated through relit mesh rendering. This distinction matters in the matched-control study: when utilization and distortion freedom are denied, the PBR gain disappears.

\paragraph*{Relightable appearance and Gaussian/surfel representations.}
Our PBR evaluation sits on standard microfacet and production-shading foundations, including Cook--Torrance reflectance \cite{cook1982_reflectance_model}, GGX-style microfacet distributions \cite{walter2007_microfacet_refraction}, Disney's production BRDF \cite{burley2012_disney_brdf}, and Unreal Engine 4's real-time PBR model \cite{karis2013_real_shading_ue4}. The differentiable-rendering side is related to OpenDR \cite{loper2014_opendr}, Neural 3D Mesh Renderer \cite{kato2018_neural_3d_mesh_renderer}, Soft Rasterizer \cite{liu2019_soft_rasterizer}, DIB-R \cite{chen2019_dibr}, differentiable Monte Carlo edge sampling \cite{li2018_differentiable_monte_carlo}, nvdiffrast \cite{laine2020_nvdiffrast_modular_primitives}, and Mitsuba 2 \cite{nimier_david2019_mitsuba2}. Relightable 3D representations have recently made strong progress in Gaussian and surfel pipelines, including scattering/shadow decomposition \cite{zheng2026_ssd_gs_scattering_shadow}, dynamic relighting cues \cite{kaleta2026_lumimotion_improving_gaussian_relighting}, appearance-disentangled feed-forward Gaussian splatting \cite{herau2026_spectralsplat_appearance_disentangled_feed}, high-frequency Gaussian/surfel reconstruction \cite{watanabe2026_neural_gabor_splatting_enhanced,dai2026_surfelsplat_learning_efficient_generalizable,gomez2026_blobs_spokes_high_fidelity}, and generated-asset systems such as GET3D, DreamFusion, and Trellis \cite{gao2022_get3d_generative_model,poole2023_dreamfusion_text_to_3d,xiang2024_trellis_structured_3d_latents}. These methods motivate using relighting as an evaluation signal, but they often bypass the mesh-atlas bottleneck or optimize scene representations directly. Our contribution is narrower: we retain a standard mesh, a standard single UV atlas, and a fixed texture budget, then use PBR relighting residuals to decide whether a local atlas repair should be deployed.
